\chapter{Crying Spells}

\section*{Definition}
Episodes of uncontrollable or excessive tearfulness that occur unexpectedly or with minimal provocation, often unrelated to the severity of the triggering event.

\section*{When to See a Doctor}
See your doctor if:
\begin{itemize}
  \item Crying spells with suicidal thoughts, hopelessness, sleep or appetite changes, or lasting more than 2 weeks with functional impairment
\end{itemize}

\section*{How It Develops}
Crying spells can reflect mood disorders, grief, stress, hormonal changes, medication effects, or neurologic conditions. The cause cannot be determined from tearfulness alone and depends on timing, associated symptoms, and functional impact.

\section*{Causes}
\begin{itemize}
  \item Perimenopausal and menopausal hormonal fluctuations
  \item Premenstrual syndrome or premenstrual dysphoric disorder
  \item Major depressive disorder
  \item Anxiety disorders (generalised, panic)
  \item Adjustment disorder with depressed mood
  \item Post-traumatic stress disorder
  \item Pseudobulbar affect (neurological condition)
  \item Chronic stress and burnout
  \item Bereavement or grief
  \item Medication side effects (hormonal contraceptives, corticosteroids)
\end{itemize}

\section*{Related Symptoms}
\begin{itemize}
  \item Depression
  \item Anxiety
  \item Mood swings
  \item Irritability
  \item Fatigue
\end{itemize}


\section*{Medical Evaluation}
\begin{itemize}
  \item Your doctor will take a detailed history of your mood, stress levels, life circumstances, and crying patterns.
  \item Screening questionnaires for depression, anxiety, and perimenopausal symptoms may be used.
   \item Targeted blood tests, such as thyroid testing or a blood count, may be ordered when the history or examination suggests a medical contributor; routine sex-hormone testing is not usually needed to diagnose perimenopause.
  \item A mental health assessment will help differentiate primary psychiatric conditions from hormonal causes.
  \item Referral to a psychiatrist, psychologist, or endocrinologist may be arranged if indicated.
\end{itemize}

\section*{Treatment Options}
\begin{itemize}
  \item Counseling or psychotherapy (cognitive-behavioral therapy, interpersonal therapy) is first-line for mood-related crying spells.
   \item Menopausal hormone therapy may be considered for bothersome vasomotor symptoms after an individualized discussion of benefits and risks; it is not a general treatment for tearfulness.
  \item Antidepressants such as SSRIs or SNRIs may be prescribed for depression or anxiety.
  \item Lifestyle modifications including improved sleep, nutrition, and stress management are essential components.
  \item Regular follow-up ensures treatment effectiveness and allows timely adjustments.
\end{itemize}

\section*{Self-Care and Home Management}
\begin{itemize}
  \item Practice stress-reduction techniques such as deep breathing, meditation, or gentle yoga.
  \item Maintain a consistent sleep schedule and prioritize restful sleep.
  \item Engage in regular physical activity to support mood regulation.
  \item Keep a mood diary to identify triggers and patterns in crying episodes.
  \item Reach out to trusted friends, family, or a support group for emotional support.
\end{itemize}

\section*{References}
\begin{thebibliography}{99}
\bibitem{crying_spells:ref1} Lapid MI, Rummans TA. ``Pseudobulbar affect: a review of current understanding and management.'' J Clin Psychiatry. 2009;70(8):1119-1125.
\bibitem{crying_spells:ref2} Work SS, Colamonico JA, Bradley WG, et al. ``Pseudobulbar affect: an under-recognized and under-treated neurological disorder.'' Adv Ther. 2011;28(7):586-601.
\bibitem{crying_spells:ref3} Alexopoulos GS. ``Depression in the elderly.'' Lancet. 2005;365(9475):1961-1970.
\bibitem{crying_spells:ref4} American Psychiatric Association. ``Major depressive disorder.'' In: Diagnostic and Statistical Manual of Mental Disorders. 5th ed, Text Revision. APA; 2022.
\bibitem{crying_spells:ref5} Cummings JL. ``Emotional lability and pseudobulbar affect.'' In: Neuropsychiatry and Behavioral Neuroscience. Oxford University Press; 2021.
\bibitem{crying_spells:ref6} Tampi RR. ``Pseudobulbar affect in adults: treatment.'' UpToDate. Wolters Kluwer; 2025.
\end{thebibliography}
